\begin{abstract}
How reliably can black-box behavioral probing distinguish truthful from deceptive language model outputs? We evaluate this question under controlled conditions across seven models (3B--70B, three families) using three reusable methodological controls, and find that behavioral detection primarily captures instruction-following artifacts. (1)~A \emph{prompt-equalized control}---identical neutral prompts, behavioral differences from knowledge conflict only---drops classification accuracy from 93--100\% to 52--69\% under independent cross-family extraction. (2)~\emph{Cross-family feature extraction} reveals that extractor choice affects accuracy by 9--16\,pp across three target models (reflecting both same-family bias and extractor capability differences); the commonly reported 61--84\% range is an upper bound; the scale trend holds under 2 of 3 extractors but is extractor-sensitive at the 7B point. (3)~A \emph{hedging-word regex baseline} matches or exceeds the full LLM pipeline under equalization, confirming the residual signal is surface-level. The equalized signal exhibits a two-tier structure: models $\leq$7B cluster at 61--74\% while models $\geq$14B cluster at 82.5--84\% ($p < 0.0001$); a within-family control (Qwen 7B: 74\% vs.\ Qwen 14B: 82.5\%) confirms this structure holds within a single model family. No individual adjacent increment is statistically significant (Holm-Bonferroni corrected). An instructed-matched control on five models confirms instruction-following as the dominant confound (+7.5--31\,pp across five models). Multi-turn interrogation is marginal under instructed conditions (+1--3\,pp) but essential under equalization, where the single-question-to-adaptive gap scales from 0\,pp at 3B to +21.5\,pp at 70B. We study instructed roleplay deception, not autonomous strategic deception; these results establish evaluation controls that future behavioral deception detection work should adopt.
\end{abstract}
